\section{A tight output or a descent at a probe}\label{sec:vertex}

The output floors transmit information down a subtree and back to
its root. Fix a probe \(x\) for bucket \(P\), and at each module use
the local coordinates of the block containing \(x\). Under a coordinate
restriction, a floor is \emph{retained} when its row is supported on
the retained coordinates. At the current module one may raise each
retained output floor independently from zero to a prescribed value.
This includes floors on propagated outputs. The descendant systems
initially keep their original inequalities.

Call a module \emph{ready} if the restriction keeps its coordinates
associated to \(W_\ell,p_\ell,F_{\ell,P},G_{\ell,P}\) for every label
\(\ell\). At height zero it suffices to keep the output coordinate for
\(P\). Readiness is a condition on this one module.

\begin{lemma}[Descent at a probe]\label{lem:descent}
The parameters at each height can be chosen with a positive allowance
for independent raises of the module's output floors. For every choice
of raises within that allowance, each vertex of the restricted slice
of a ready subtree has at least one of the following properties.
\begin{enumerate}
\item A quick output of the module in bucket $P$ is at its floor.
\item The module has positive height and has a label $\ell$ with the
following property for every child at $\ell$. After allowed raises of
that child's output floors, the point restricted to its subtree is a
vertex. A retained child output is at its prescribed floor precisely
when its propagated output is at its floor in the parent.
\end{enumerate}
The second alternative includes a label with no children.
\end{lemma}

The two quick outputs in bucket \(P\) enforce this alternative in
different ways. If neither kind is tight, we first find a bundle of
labels on which the link between \(W\) and the price \(p\) is slack.
The \(F\)-outputs are responsible for this freedom. The tournament
then singles out a label in the bundle, and a \(G\)-output forces its
price into the interval in which the parent floors become admissible
raises in every child.

\begin{proof}
At height zero take allowance $1/2$. At the probe for $P$, the
absolute-value bound on its output has value $-1$, and every other
present output is at least $2$. All pair bounds touching the $P$
output are therefore strict. Unless that output is at its floor, it
can move in both directions.

For a module of positive height we prove that \(\varepsilon\lambda/4\)
is an allowance. The child systems and their shifts are already fixed.
Suppose there are arbitrarily small \(\sigma>0\) for which the
assertion fails. Take a sequence \(\sigma\downarrow0\) and, for each
term, a counterexample vertex with prescribed raises
\(0\le a_j\le\varepsilon\lambda/4\). Pass to a subsequence fixing the
module, probe and coordinate restriction. In particular, every
retained target quick output is strictly above its prescribed floor.

At this sample we need only the current block's rows at the module
(Section~\ref{sec:layers}). By readiness the triangular change
\eqref{eq:module-coordinates} is invertible on the retained variables.
All prices and target outputs are present. An output from another
bucket is used only when its associated coordinate is retained.

\smallskip
\noindent\emph{Bounds for the children.}
Keep the current module fixed. Each child point is then a vertex of its
own restricted system augmented by the propagated floors and pair rows
from the parent. The frozen section is a product
over children, since a pair of propagated outputs belongs to a single
class from a single child. Dividing these extra rows by $\varepsilon$
gives fixed normals and right-hand sides $a_j/\varepsilon-p_\ell/\varepsilon$
and $\delta-2p_\ell/\varepsilon$, respectively. The price floor
$p_\ell\ge-K_0\varepsilon$ and the bound $a_j\le\varepsilon\lambda/4$
bound these right-hand sides above. For a tight extra row its right-hand
side also has a lower bound. Indeed it then equals the corresponding
physical child output or the sum of two such outputs, and a retained
parent row includes the supports of the relevant child floors. Indeed
the output expressions have nonnegative coefficients in the coordinates
\eqref{eq:module-coordinates}, so their supports cannot cancel in a sum.
The child floors give fixed lower bounds at \(x\). A vertex is determined by
linearly independent tight rows of full rank in the retained child
coordinates. Their normals come from finitely many fixed matrices, and
all their tight right-hand sides have just been bounded on both sides.
Thus the child coordinates remain bounded as \(\sigma\downarrow0\).

\smallskip
\noindent\emph{A common limiting state.}
Now freeze the children instead. The current node is a vertex of this
section. Recall from \eqref{eq:layer-state} that $z=W-\sigma g(x)\mathbf 1$
and that its coarse constraint is $z^+\in D Z_x$, with
$D=1+M_0\sum_\ell z_\ell^-$. Hold this particular value of $D$ fixed
and replace $W$ and all its constraints by
\[
  \widehat z\in D Z_x,\qquad
  F_{\ell,Q}\ge\widehat z_\ell-\sigma o_Q(x),\qquad
  p_\ell\ge\widehat z_\ell-3\mu .
\]
The point with $\widehat z=z^+$ is feasible. At negative entries of
$z$, the displayed output links follow from $F_{\ell,Q}\ge0$ and
$p_\ell\ge-K_0\varepsilon>-3\mu$ for small $\sigma$.
It is still a vertex. Any two-sided motion in the new section fixes
the zero entries of $\widehat z$ by nonnegativity. Restore their
original $z$-values and let the other entries move with $\widehat z$
on a sufficiently short segment. All negative parts stay fixed, so the
original dilation is preserved. The original links follow from
$z\le\widehat z$, and $W\ge0$ holds since moving entries of $z$
are positive. This would lift any nonzero two-sided motion to the
original section.

Divide the current variables by \(D\), writing dots for the
result. Since \(W\ge0\),
\[
        1\le D\le 1+M_0V\sigma g(x).
\]
In the divided coordinates every constraint can be written with a
normal independent of \(\sigma\). For the tournament constraints,
first multiply \eqref{eq:module-tournament} by \(\varepsilon\).
The frozen child coordinates are bounded, so propagated floors and
pair rows give lower bounds on \(\dot p_\ell\) of order
\(O(\varepsilon)\). The limiting inequalities reduce to
\[
 \dot{\widehat z}\in Z_x,\qquad
 \dot p_\ell\ge\dot{\widehat z}_\ell,\qquad
 \dot F_{\ell,Q}\ge\dot{\widehat z}_\ell,\qquad
 \dot G_{\ell,Q}\ge\dot p_\ell .
\]
All remaining rows follow from these inequalities. Their right-hand
sides tend to zero, and their left-hand sides are nonnegative
combinations of the displayed variables. Take a further subsequence
on which one nonsingular matrix of tight rows determines the divided
vertex. Its right-hand sides converge, so its solution converges and
the same independent rows remain tight in the limit. The limit is
therefore a vertex of the displayed polyhedron. The differences
$\dot p-\dot{\widehat z}$, $\dot F-\dot{\widehat z}$ and $\dot G-\dot p$,
taken label by label, identify the limiting system
with $Z_x$ times an orthant. At a vertex all those differences vanish and
$\dot{\widehat z}$ equals a vertex $z_*$ of $Z_x$. Since $Z_x$ is
the convex hull of the sample states plus the nonnegative orthant,
$z_*$ is one of the states.

\smallskip
\noindent\emph{Freedom of the prices at zero labels.}
We claim that whenever $z_{*,\ell}=0$ the original price satisfies, for all sufficiently small $\sigma$ in our sequence,
\begin{equation}\label{eq:zero-price-link}
                  p_\ell>z_\ell-3\mu .
\end{equation}
Call a target $F$ or $G$ \emph{paired} if one of its pair inequalities
is tight. Pass to a fixed pattern of pairings for the target $F$'s.
For a paired target $F$, the other-bucket partner is at most $\mu$
by the tight pair bound and nonnegativity. The partner's link then gives
$\widehat z_\ell\le\mu$. Thus the target link to $\widehat z_\ell$ is
strict, as $\widehat z_\ell-\sigma<0$. The price floor
$p_\ell\ge-K_0\varepsilon>-2\mu$ implies
\eqref{eq:zero-price-link} at such a label.

It remains to show that a target \(F_{\ell,P}\) at a zero label
cannot be unpaired. In the divided system the target link has an
offset of order \(\sigma\), whereas the other tight rows have offsets
of order at most \(\mu=\sigma^2\). For each unpaired target \(F\),
use \((F_{\ell,P}+\sigma)/D\) in place of \(F_{\ell,P}/D\).
Its floor and pair rows are slack, so the only tight row involving
this variable is its link to \(\widehat z_\ell\). The translation
cancels the offset in that row. For paired targets the link is already
slack. In the translated tight system the right-hand sides therefore differ
from their limits by \(O(\mu)\). The rows of \(Z_x\) are fixed,
the other-bucket \(F\)-links have zero offset, and the remaining price,
pair and floor terms have size \(O(\mu)\) or \(O(\varepsilon)\).
The tight basis has a fixed inverse, which gives the same error bound
for each translated variable. At a zero label the limiting
target output is zero, while any unpaired translated output is at least
$\sigma/D$. This contradicts $\mu=\sigma^2$ and $D\to1$,
establishing \eqref{eq:zero-price-link}.

\smallskip
\noindent\emph{One label for all children.}
Let $s$ be the site containing the probe in its current block $E$.
Each tournament tuple for another site $s'$ has a label with
$z_{*,\ell}>0$, by the bound of $r$ zeros per foreign bundle. At
that label $p_\ell$ tends to a positive value, whence
$u_\ell^{s'}\to-\infty$. All other terms are bounded above by $K_1$
before weighting, so every foreign-site tournament row is strict.

Choose a bundle $H_i(s)$ on which $z_*$ vanishes. Fewer than $2^r$
of its positions have $u_\ell^s>0$. Otherwise there is a transitive
$(r+1)$-tuple among them, contradicting its tournament row. Indeed
any tournament on $2^r$ vertices has a vertex beating at least
$2^{r-1}$ others, and induction supplies the rest of the tuple. Choose
a position outside this positive set beating every position in it,
using the defining property of the tournament on positions. Pass to a
subsequence so its label $\ell\in H_i(s)$ is fixed. We have $u_\ell^s\le0$.
If $u_\ell^s<-\delta/2$, every current-site tournament row containing
this label is strict. Indeed the label determines $i$, so every
such tuple lies in $H_i(s)$. All positive terms follow
$\ell$ in the order and each has weight at most $1/M_1$ times that
of $\ell$. Their combined contribution, relative to its weight, is
at most $rK_1/M_1<\delta/2$.

Express the selected price in local child units by
\[
  \theta=-p_\ell/\varepsilon-D_{E_s}(x)
        =\delta+u_\ell^s-R_{t-1}(x)\le\delta .
\]
A propagated output \eqref{eq:module-propagation} is
$\varepsilon(v_j-\theta)$ for the local child output $v_j$.
The child pair sum exceeds $10\delta$, so a
retained pair has parent sum $\varepsilon(v_j+v_k-2\theta)>8\varepsilon\delta$,
strictly above its bound $\varepsilon\delta$.
At this label the target $G$ link has bound $\varepsilon(-\theta-\lambda)$,
and each other-bucket $G$ link has bound $\varepsilon(-\theta+\lambda)$.

If $G_{\ell,P}$ is unpaired, it must attain its link bound, since its floor is slack and it could otherwise move alone. It follows that
$\theta<-\lambda$. We can then move the price $p_\ell$ and every
present $G$ at this label by the same small amount in either
direction. Indeed,
$u_\ell^s=\theta-\delta+R_{t-1}(x)<-\delta/2$, using
$R_{t-1}(x)\le\delta/8$. All tournament rows at this price are
therefore slack. Its price floor is slack too, since at that floor
$u_\ell^s=K_0-T_{E_s}(x)-\delta>0$.
Its link from $W$ is slack by \eqref{eq:zero-price-link}.
All present other-bucket $G$'s exceed $2\varepsilon\lambda$, so their
floors and pair rows are strict. The target $G$ has slack floor and
no tight pair. Each retained propagated output also exceeds \(\varepsilon\lambda\),
since \(v_j\ge0\), while its floor is at most
\(\varepsilon\lambda/4\). Its pair rows were shown strict above.
The price motion preserves the \(G\)-links, contradicting extremality.

Thus the target $G$ is paired. Its other-bucket partner has value at most
$\varepsilon\lambda/4$. Substituting the other-bucket link gives
$3\lambda/4\le\theta\le\delta$. Freeze the selected price. For a
propagated output with floor $a_j\in[0,\varepsilon\lambda/4]$ the
parent now requires
\[
                  v_j\ge \theta+a_j/\varepsilon .
\]
The right-hand side is positive and at most $\delta+\lambda/4\le2\delta$,
within the child's floor allowance. Because the price coordinates are
present, this propagated row survives precisely when the child output
floor survives. All other parent rows on these child coordinates
are the slack propagated pair bounds. A two-sided motion of any one
child point respecting the new floors would consequently lift to a
two-sided motion of the original vertex. Each child point is a vertex
with these raises, and equality in a raised floor gives equality in
the parent's propagated floor. This proves the second alternative for
our common label, contradicting the choice of the sequence.

The contradiction excludes a sequence of failures for a fixed module,
probe and restriction, even when the raises vary along the sequence.
There are finitely many such triples at this height, so one sufficiently
small \(\sigma\) works for all of them. This choice uses only the local
rows at the current height and the systems and shifts already fixed at
lower heights. Once the floor allowance is known, the current shifts
can be chosen as in Section~\ref{sec:layers}.
\end{proof}
